\section{Frontend Technologies}
In this section we will explain the technologies used and why we chose them in the frontend of the project.

\subsection{Shadcn/UI}
Shadcn/UI is an open-source library of pre-built components based on Radix UI and Tailwind CSS. To use them we just need to add the components to the project and use them in the code. 
Tailwind CSS is a CSS framework that allows us to style the components with a single line of code, making the development process faster and more efficient.
In our webproject we used the components to agilize the development process and to have a more consistent design throughout the app. Most of the components are already the way we wanted them to be, but we can and did customize some of them to fit our needs using Tailwind CSS and changing the default styles.

\subsection{React}
React is a JavaScript library for building user interfaces. Helps us create reusable components and manage the state of the app, while keeping the code clean and easy to understand.
It uses a Virtual DOM which is a lightweight representation of the DOM. When the state of the app changes, React updates the Virtual DOM and compares it to the previous Virtual DOM to determine the changes needed to the actual DOM. This is more efficient than updating the DOM directly.

In our project we used React to build the frontend of the webapp, from the pages layouts to the components. React allows us to update the UI allongside the data, making the app more fluid and responsive. One of the usecases of React is the use of hooks, which are functions that allow us to manage the state of the app like \texttt{useState}, \texttt{useEffect}, etc. This is very useful when we need to manage the state of the app, like in the case of the form validation.

Additionally, React provides access to a lot of third party libraries that can help us with the development of the app.

\subsection{Next.js}
Next.js is an open-source React framework developed by Vercel. It is a React framework that provides a set of features like server-side rendering, static site generation, and routing. Next.js is known for its performance and simplifing routing.

In our project we chose Next.js as the framework for the frontend because it provides a lot of features that are useful for the development of the app and the use of React. The built-in routing system allowed us to create pages only by adding a new file in the \texttt{app} folder, with no need to configure the routes in a different file.

Overall, Next.js is a powerful framework that provides a lot of features that are useful for the development of the app and the use of React.

\subsection{React Native}
React Native is a open-source framework, developed by Meta Platforms, used to build mobile applications, to both Android and iOS, using React, as well as web and desktop applications.
It allows us to create native mobile applications using the same principles and components as React, but with the ability to access native device features.

In our project we used React Native to build the mobile app, allowing to make the app available to both Android and iOS audience.

\subsection{Expo}
Expo is an open-source platform built on top of React Native that simplifies the development process of mobile applications. It provides a set of tools and services that streamline the development, testing, and deployment of React Native apps.
Expo includes a development environment, a set of pre-built components, and a build service that allows us to create and publish mobile applications without the need for complex native configurations.

We took advantage of Expo’s managed workflow, using modules such as expo-status-bar, expo-linear-gradient, and expo-constants to simplify access to native features without manual native configuration. 
Additionally, Expo Router was used to handle navigation, and the build service allowed us to generate production-ready builds for both Android and iOS platforms efficiently.

\subsection{TypeScript}
TypeScript is a statically typed programming language that is a superset of JavaScript. It is a popular language developed by Microsoft for building large-scale applications, and it adds optional static types, classes and interfaces to JavaScript.

By adding types to the code it allows us to catch errors early instead of at runtime, making the development process more efficient and the code more robust.

In our project we used TypeScript in the frontend to improve the development experience with a less runtime errors and a better code organization. The fact that has integration with React and Next.js is a plus, as well as the fact that it is a very popular language in the industry. 

Overall, TypeScript is a powerful language that can help us have more clarity in the code, catch errors early and make the development process more efficient.

\subsection{Jest}
Jest is a JavaScript testing framework developed and maintained by Facebook (Meta), with a focus on simplicity and ease of use. This framework provides integrated features such as mocking, code coverage, and snapshot testing, eliminating the need for complex configurations.
In our project, we used Jest to perform unit tests on the components of both web and mobile applications. These tests allow us to verify that components behave as expected, both individually and when integrated together, facilitating early error detection and code maintenance.